\subsection{Computational performance}\label{sec:results}

Which configuration most often certifies a solution, and does its
relative performance vary with risk objective, training size or
budget? Table~\ref{tab:new20-gaps-times} reports all 45 runs per
method and objective. The gaps describe optimization accuracy of the
recorded formulations; they are not differences in out-of-sample
protection. Mean observed time must be read jointly with certification
and final gap because unresolved runs reached the stopping limit.

\begin{center}
\begin{minipage}{0.98\linewidth}\centering\small
\setlength{\tabcolsep}{3pt}
\captionof{table}{Reduced \texttt{new20x20}: percentage certified to
relative MIP gap $10^{-3}$ ($100N_{\rm opt}/N$, $N=45$ per row), mean
and maximum final relative MIP gaps (\%), and mean observed solver
time (s). All 45 runs, including uncertified runs and recorded time
overruns, enter the gap and time statistics. Generated from the 810
run-level records by \texttt{analysis/reproduce\_results.py}.}
\label{tab:new20-gaps-times}
\begin{tabular}{llrrrrr}
\toprule
Objective & Method & $N_{\rm opt}/N$ & Mean gap & Max gap & Mean time & SD time \\
\midrule
Expected & Extensive SAA & 42/45 & 0.75 & 21.79 & 409.8 & 545.8 \\
Expected & LP B\&B + Root & 26/45 & 4.44 & 15.71 & 878.6 & 843.0 \\
Expected & LP B\&B + Standard & 26/45 & 4.65 & 16.53 & 890.8 & 838.4 \\
Expected & LP B\&B + Coverage-exp & 34/45 & 1.43 & 12.52 & 577.9 & 788.4 \\
Expected & LP B\&B + Path-exp & 39/45 & 0.95 & 10.78 & 356.5 & 632.3 \\
Expected & Combinatorial B\&B & 43/45 & 0.18 & 4.00 & 212.4 & 427.1 \\
\midrule
CVaR & Extensive SAA & 6/45 & 36.50 & 100.00 & 1667.3 & 374.8 \\
CVaR & LP B\&B + Root & 15/45 & 11.18 & 26.60 & 1218.3 & 834.2 \\
CVaR & LP B\&B + Standard & 16/45 & 10.74 & 28.19 & 1209.4 & 843.8 \\
CVaR & LP B\&B + Coverage-exp & 19/45 & 5.06 & 15.90 & 1101.0 & 854.4 \\
CVaR & LP B\&B + Path-exp & 21/45 & 3.51 & 12.40 & 1003.5 & 872.4 \\
CVaR & Combinatorial B\&B & 19/45 & 5.31 & 21.67 & 1254.4 & 784.8 \\
\midrule
Mean--CVaR & Extensive SAA & 13/45 & 17.76 & 99.97 & 1456.9 & 598.2 \\
Mean--CVaR & LP B\&B + Root & 15/45 & 8.75 & 22.41 & 1217.2 & 836.4 \\
Mean--CVaR & LP B\&B + Standard & 17/45 & 9.38 & 23.22 & 1207.4 & 840.4 \\
Mean--CVaR & LP B\&B + Coverage-exp & 22/45 & 3.46 & 13.21 & 1051.5 & 849.2 \\
Mean--CVaR & LP B\&B + Path-exp & 30/45 & 2.10 & 10.86 & 813.1 & 824.9 \\
Mean--CVaR & Combinatorial B\&B & 21/45 & 4.24 & 16.29 & 1182.4 & 800.9 \\
\bottomrule
\end{tabular}

\end{minipage}
\end{center}

For expected loss, Combinatorial B\&B certifies 43/45 runs, has the
lowest mean gap (0.18\%) and mean observed time (212.4\,s), and its
maximum final gap is 4.00\%. Extensive SAA certifies 42/45 with mean
time 409.8\,s; Path-exp certifies 39/45 with 356.5\,s. Thus
Combinatorial B\&B is the strongest recorded expected-loss
configuration across these joint criteria. The risk objectives change
the ordering. For CVaR, Path-exp certifies 21/45, compared with 19/45
for each of Coverage-exp and Combinatorial B\&B. Path-exp also has the
smallest mean gap (3.51\%), smallest maximum gap (12.40\%) and lowest
mean observed time (1,003.5\,s). For mean--CVaR, Path-exp certifies
30/45 versus 22/45 for Coverage-exp and 21/45 for Combinatorial
B\&B, again with the lowest mean and maximum gaps (2.10\% and
10.86\%) and mean observed time (813.1\,s). These are finite-panel
computational comparisons, not statements of decision quality.

\begin{figure}[t]
\centering
\begin{tikzpicture}
\begin{groupplot}[group style={group size=3 by 1,horizontal sep=9pt},
  width=0.315\linewidth,height=0.275\linewidth,xmin=0,xmax=1800,
  ymin=0,ymax=1,xtick={0,600,1200,1800},ytick={0,0.2,0.4,0.6,0.8,1},
  xlabel={Solver time (s)},xlabel style={font=\scriptsize},
  tick label style={font=\tiny},title style={font=\small},
  grid=major,grid style={black!15},enlarge x limits=false]
\nextgroupplot[title={Expected},ylabel={Fraction certified / 45},ylabel style={font=\scriptsize},legend to name=certlegend,legend columns=3,legend style={draw=none,font=\scriptsize}]
\addplot+[const plot,mark=none,thin,color=gray!80!black] coordinates {(0.0000,0.0000000) (12.1499,0.0222222) (15.5780,0.0444444) (19.2062,0.0666667) (19.7313,0.0888889) (22.3492,0.1111111) (22.7488,0.1333333) (24.2890,0.1555556) (36.6375,0.1777778) (38.9064,0.2000000) (41.8699,0.2222222) (45.2682,0.2444444) (48.2516,0.2666667) (55.1546,0.2888889) (57.1705,0.3111111) (64.6277,0.3333333) (65.8659,0.3555556) (78.4806,0.3777778) (86.0197,0.4000000) (90.7905,0.4222222) (97.3513,0.4444444) (109.5172,0.4666667) (116.7803,0.4888889) (119.9584,0.5111111) (127.2888,0.5333333) (140.5663,0.5555556) (144.5641,0.5777778) (168.4127,0.6000000) (193.3160,0.6222222) (248.8954,0.6444444) (403.9395,0.6666667) (407.2200,0.6888889) (498.9915,0.7111111) (545.4261,0.7333333) (558.7232,0.7555556) (566.8024,0.7777778) (589.5554,0.8000000) (910.6997,0.8222222) (998.6312,0.8444444) (1059.1241,0.8666667) (1358.5070,0.8888889) (1367.1474,0.9111111) (1462.1231,0.9333333) (1800.0000,0.9333333)};
\addlegendentry{SAA}
\addplot+[const plot,mark=none,thin,color=red!75!black] coordinates {(0.0000,0.0000000) (4.6272,0.0222222) (4.8418,0.0444444) (5.7517,0.0666667) (8.9113,0.0888889) (9.6470,0.1111111) (10.9565,0.1333333) (14.4311,0.1555556) (16.5968,0.1777778) (28.6764,0.2000000) (29.6264,0.2222222) (29.6303,0.2444444) (30.0510,0.2666667) (34.1671,0.2888889) (41.0872,0.3111111) (44.4412,0.3333333) (44.7226,0.3555556) (57.0497,0.3777778) (126.9117,0.4000000) (300.7080,0.4222222) (330.9929,0.4444444) (348.1397,0.4666667) (358.9008,0.4888889) (468.3254,0.5111111) (596.3445,0.5333333) (680.7388,0.5555556) (1680.0052,0.5777778) (1800.0000,0.5777778)};
\addlegendentry{Root}
\addplot+[const plot,mark=none,thin,color=orange!85!black] coordinates {(0.0000,0.0000000) (3.1048,0.0222222) (4.8732,0.0444444) (6.4458,0.0666667) (6.7589,0.0888889) (8.0769,0.1111111) (8.5478,0.1333333) (9.6328,0.1555556) (13.1416,0.1777778) (13.8146,0.2000000) (21.1993,0.2222222) (25.3963,0.2444444) (25.9327,0.2666667) (31.4800,0.2888889) (39.9890,0.3111111) (55.9509,0.3333333) (56.0007,0.3555556) (72.5836,0.3777778) (156.8989,0.4000000) (218.5410,0.4222222) (226.8395,0.4444444) (259.1950,0.4666667) (606.8967,0.4888889) (657.9411,0.5111111) (731.7437,0.5333333) (1140.8392,0.5555556) (1439.8117,0.5777778) (1800.0000,0.5777778)};
\addlegendentry{Standard}
\addplot+[const plot,mark=none,thin,color=blue!75!black] coordinates {(0.0000,0.0000000) (5.5601,0.0222222) (7.0530,0.0444444) (7.0862,0.0666667) (7.1337,0.0888889) (7.5863,0.1111111) (11.3390,0.1333333) (13.4299,0.1555556) (13.7843,0.1777778) (16.8579,0.2000000) (16.9001,0.2222222) (17.1779,0.2444444) (17.6503,0.2666667) (18.1093,0.2888889) (21.1686,0.3111111) (21.6242,0.3333333) (22.0268,0.3555556) (26.2592,0.3777778) (35.0162,0.4000000) (35.6190,0.4222222) (35.6872,0.4444444) (41.7065,0.4666667) (43.5143,0.4888889) (44.0839,0.5111111) (50.1900,0.5333333) (70.0403,0.5555556) (103.3690,0.5777778) (116.1591,0.6000000) (119.9764,0.6222222) (351.8233,0.6444444) (366.9909,0.6666667) (603.8314,0.6888889) (651.5410,0.7111111) (1311.0613,0.7333333) (1607.3191,0.7555556) (1800.0000,0.7555556)};
\addlegendentry{Coverage-exp}
\addplot+[const plot,mark=none,thin,color=green!65!black] coordinates {(0.0000,0.0000000) (1.7714,0.0222222) (2.0903,0.0444444) (2.5883,0.0666667) (2.9036,0.0888889) (3.1866,0.1111111) (4.5763,0.1333333) (4.6803,0.1555556) (4.9066,0.1777778) (5.8047,0.2000000) (6.1114,0.2222222) (7.2634,0.2444444) (8.0601,0.2666667) (8.5570,0.2888889) (8.5911,0.3111111) (10.0121,0.3333333) (10.3810,0.3555556) (11.5185,0.3777778) (11.8849,0.4000000) (13.1450,0.4222222) (15.9267,0.4444444) (16.0572,0.4666667) (17.7320,0.4888889) (19.0621,0.5111111) (21.5110,0.5333333) (22.0115,0.5555556) (23.4375,0.5777778) (24.5656,0.6000000) (35.5872,0.6222222) (40.0316,0.6444444) (44.1848,0.6666667) (100.1079,0.6888889) (247.5876,0.7111111) (260.2609,0.7333333) (263.4319,0.7555556) (267.9180,0.7777778) (759.4238,0.8000000) (768.2701,0.8222222) (952.4643,0.8444444) (1187.4260,0.8666667) (1800.0000,0.8666667)};
\addlegendentry{Path-exp}
\addplot+[const plot,mark=none,thin,color=violet!75!black] coordinates {(0.0000,0.0000000) (8.0168,0.0222222) (8.0362,0.0444444) (10.6377,0.0666667) (16.3465,0.0888889) (17.7914,0.1111111) (18.3737,0.1333333) (18.6214,0.1555556) (21.2004,0.1777778) (24.1056,0.2000000) (26.4642,0.2222222) (28.3320,0.2444444) (30.5424,0.2666667) (32.4178,0.2888889) (38.1508,0.3111111) (39.5477,0.3333333) (41.1172,0.3555556) (46.2337,0.3777778) (50.3529,0.4000000) (51.6455,0.4222222) (51.6737,0.4444444) (53.3525,0.4666667) (55.6016,0.4888889) (61.3396,0.5111111) (65.5695,0.5333333) (67.6382,0.5555556) (68.5008,0.5777778) (69.1737,0.6000000) (70.9838,0.6222222) (79.6385,0.6444444) (81.2436,0.6666667) (86.0619,0.6888889) (91.0043,0.7111111) (93.3924,0.7333333) (95.4004,0.7555556) (114.3509,0.7777778) (228.3476,0.8000000) (247.4937,0.8222222) (254.9224,0.8444444) (256.3191,0.8666667) (284.8170,0.8888889) (641.2623,0.9111111) (1104.4254,0.9333333) (1175.1750,0.9555556) (1800.0000,0.9555556)};
\addlegendentry{Combinatorial}
\nextgroupplot[title={CVaR},yticklabels={}]
\addplot+[const plot,mark=none,thin,color=gray!80!black] coordinates {(0.0000,0.0000000) (226.0480,0.0222222) (574.3726,0.0444444) (586.5184,0.0666667) (807.5886,0.0888889) (1137.8618,0.1111111) (1487.2829,0.1333333) (1800.0000,0.1333333)};
\addplot+[const plot,mark=none,thin,color=red!75!black] coordinates {(0.0000,0.0000000) (3.2267,0.0222222) (3.6386,0.0444444) (9.7191,0.0666667) (10.7087,0.0888889) (16.1555,0.1111111) (18.3812,0.1333333) (24.4677,0.1555556) (43.4074,0.1777778) (50.3609,0.2000000) (56.1632,0.2222222) (56.4050,0.2444444) (93.9259,0.2666667) (96.4051,0.2888889) (111.3530,0.3111111) (193.3495,0.3333333) (1800.0000,0.3333333)};
\addplot+[const plot,mark=none,thin,color=orange!85!black] coordinates {(0.0000,0.0000000) (2.7316,0.0222222) (4.4185,0.0444444) (9.0208,0.0666667) (10.4435,0.0888889) (11.1994,0.1111111) (16.8819,0.1333333) (18.6078,0.1555556) (25.2945,0.1777778) (27.2636,0.2000000) (27.7552,0.2222222) (33.2702,0.2444444) (43.3607,0.2666667) (53.4409,0.2888889) (73.7713,0.3111111) (87.9832,0.3333333) (1713.2787,0.3555556) (1800.0000,0.3555556)};
\addplot+[const plot,mark=none,thin,color=blue!75!black] coordinates {(0.0000,0.0000000) (6.9180,0.0222222) (7.5397,0.0444444) (11.0619,0.0666667) (12.3901,0.0888889) (14.5376,0.1111111) (20.9483,0.1333333) (28.6698,0.1555556) (38.1773,0.1777778) (40.6758,0.2000000) (53.1363,0.2222222) (55.8530,0.2444444) (58.6165,0.2666667) (65.4648,0.2888889) (89.8576,0.3111111) (90.3613,0.3333333) (309.0253,0.3555556) (356.8258,0.3777778) (508.4446,0.4000000) (525.4998,0.4222222) (1800.0000,0.4222222)};
\addplot+[const plot,mark=none,thin,color=green!65!black] coordinates {(0.0000,0.0000000) (4.2390,0.0222222) (4.9103,0.0444444) (6.5463,0.0666667) (6.7701,0.0888889) (8.9445,0.1111111) (11.3314,0.1333333) (13.7324,0.1555556) (16.3980,0.1777778) (17.1492,0.2000000) (29.9777,0.2222222) (31.3250,0.2444444) (33.1557,0.2666667) (37.8050,0.2888889) (38.3536,0.3111111) (50.2052,0.3333333) (80.4604,0.3555556) (142.8856,0.3777778) (163.0846,0.4000000) (198.4675,0.4222222) (275.4645,0.4444444) (669.5135,0.4666667) (1800.0000,0.4666667)};
\addplot+[const plot,mark=none,thin,color=violet!75!black] coordinates {(0.0000,0.0000000) (16.3939,0.0222222) (23.7880,0.0444444) (54.7128,0.0666667) (57.4902,0.0888889) (83.9574,0.1111111) (100.3576,0.1333333) (128.5330,0.1555556) (131.9417,0.1777778) (141.8217,0.2000000) (166.4643,0.2222222) (202.3615,0.2444444) (275.0289,0.2666667) (415.8665,0.2888889) (548.4473,0.3111111) (757.0346,0.3333333) (829.5376,0.3555556) (1071.2758,0.3777778) (1476.4701,0.4000000) (1739.4829,0.4222222) (1800.0000,0.4222222)};
\nextgroupplot[title={Mean--CVaR},yticklabels={}]
\addplot+[const plot,mark=none,thin,color=gray!80!black] coordinates {(0.0000,0.0000000) (96.1357,0.0222222) (117.3680,0.0444444) (131.0836,0.0666667) (277.8783,0.0888889) (385.0860,0.1111111) (415.7102,0.1333333) (583.1859,0.1555556) (634.0663,0.1777778) (644.2765,0.2000000) (910.0856,0.2222222) (917.4345,0.2444444) (1065.0141,0.2666667) (1777.6341,0.2888889) (1800.0000,0.2888889)};
\addplot+[const plot,mark=none,thin,color=red!75!black] coordinates {(0.0000,0.0000000) (3.6893,0.0222222) (4.2674,0.0444444) (5.9672,0.0666667) (7.9260,0.0888889) (13.5645,0.1111111) (14.6395,0.1333333) (18.8433,0.1555556) (40.2280,0.1777778) (50.0700,0.2000000) (52.6756,0.2222222) (59.3770,0.2444444) (70.8308,0.2666667) (91.1618,0.2888889) (121.1073,0.3111111) (167.6175,0.3333333) (1800.0000,0.3333333)};
\addplot+[const plot,mark=none,thin,color=orange!85!black] coordinates {(0.0000,0.0000000) (2.8260,0.0222222) (5.3630,0.0444444) (8.3379,0.0666667) (8.7464,0.0888889) (12.9020,0.1111111) (13.5343,0.1333333) (16.7590,0.1555556) (25.2974,0.1777778) (29.8121,0.2000000) (44.8775,0.2222222) (45.9212,0.2444444) (45.9213,0.2666667) (46.1531,0.2888889) (48.7359,0.3111111) (140.6459,0.3333333) (1635.7842,0.3555556) (1729.2611,0.3777778) (1800.0000,0.3777778)};
\addplot+[const plot,mark=none,thin,color=blue!75!black] coordinates {(0.0000,0.0000000) (7.7512,0.0222222) (8.1393,0.0444444) (8.6878,0.0666667) (9.0544,0.0888889) (22.3309,0.1111111) (22.9108,0.1333333) (24.8121,0.1555556) (28.4992,0.1777778) (39.3420,0.2000000) (46.3715,0.2222222) (50.2076,0.2444444) (59.1615,0.2666667) (90.7963,0.2888889) (113.8195,0.3111111) (115.0091,0.3333333) (159.4797,0.3555556) (167.0649,0.3777778) (178.0457,0.4000000) (701.2159,0.4222222) (806.0503,0.4444444) (1311.1031,0.4666667) (1376.2803,0.4888889) (1800.0000,0.4888889)};
\addplot+[const plot,mark=none,thin,color=green!65!black] coordinates {(0.0000,0.0000000) (3.2281,0.0222222) (3.3747,0.0444444) (4.5480,0.0666667) (8.7020,0.0888889) (9.6517,0.1111111) (10.8346,0.1333333) (14.8577,0.1555556) (15.4612,0.1777778) (18.7363,0.2000000) (26.0396,0.2222222) (32.4267,0.2444444) (33.0762,0.2666667) (36.6979,0.2888889) (37.2701,0.3111111) (61.5585,0.3333333) (70.4590,0.3555556) (74.0340,0.3777778) (87.3005,0.4000000) (94.4950,0.4222222) (108.6742,0.4444444) (127.0175,0.4666667) (135.2089,0.4888889) (148.3323,0.5111111) (649.9938,0.5333333) (836.8797,0.5555556) (1164.9157,0.5777778) (1199.2989,0.6000000) (1378.3715,0.6222222) (1531.4339,0.6444444) (1558.1323,0.6666667) (1800.0000,0.6666667)};
\addplot+[const plot,mark=none,thin,color=violet!75!black] coordinates {(0.0000,0.0000000) (23.2145,0.0222222) (31.9821,0.0444444) (44.8312,0.0666667) (62.3962,0.0888889) (74.4779,0.1111111) (85.3746,0.1333333) (88.8035,0.1555556) (96.9093,0.1777778) (143.1294,0.2000000) (154.8382,0.2222222) (160.5326,0.2444444) (221.3497,0.2666667) (251.7200,0.2888889) (410.6170,0.3111111) (536.0458,0.3333333) (549.8967,0.3555556) (860.1347,0.3777778) (873.9971,0.4000000) (1106.2123,0.4222222) (1217.6393,0.4444444) (1227.2492,0.4666667) (1800.0000,0.4666667)};
\end{groupplot}
\end{tikzpicture}
\par\smallskip
\pgfplotslegendfromname{certlegend}

\caption{Cumulative certification against observed solver time, one
panel per objective. Each method has 45 runs in the denominator at
every time; unresolved runs never enter its curve. This is a
cumulative certification profile under the 1,800\,s stopping rule,
reconstructed by \texttt{analysis/reproduce\_results.py}.}
\label{fig:new20-certified}
\end{figure}

Increasing $n$ makes the risk-objective instances harder for the best
recorded configuration. For Path-exp under CVaR, the certified counts
at $n=100,200,400$ are 9, 7 and 5 of 15; its mean final gap rises
from 2.18\% to 2.93\% and 5.43\%, and mean observed time from
753.8 to 1,032.9 and 1,223.8\,s. Under mean--CVaR the corresponding
counts are 14, 11 and 5, the mean gaps 0.28\%, 1.08\% and 4.95\%,
and observed times 450.0, 769.6 and 1,219.8\,s. Expected loss
shows a method-specific exception: Combinatorial B\&B certifies
14, 14 and 15 of 15 as $n$ grows, with mean final gaps 0.22\%,
0.27\% and 0.03\%, although its observed mean time rises from
169.4 to 177.5 and 290.5\,s. Extensive SAA instead certifies
15, 15 and 12 of 15, and its mean time increases from 35.0 to
155.5 and 1,038.8\,s. These pooled size comparisons do not isolate
the effect of $n$ from interactions with the budget.

Increasing $\alpha$ is associated with fewer certifications and larger
gaps for Path-exp on both risk objectives. For CVaR, its counts across
$\alpha=0.01,0.02,0.03$ are 15, 6 and 0 of 15, mean gaps 0.07\%,
3.07\% and 7.40\%, and mean observed times 29.0, 1,176.5 and
1,804.9\,s. For mean--CVaR the corresponding values are 15, 10
and 5, 0.06\%, 1.84\% and 4.41\%, and 22.8, 884.5 and
1,532.0\,s. Expected-loss Combinatorial B\&B certifies 15, 15
and 13 of 15; its mean gap changes from 0.01\% and 0.01\% to
0.51\%, while observed time rises from 33.5 to 72.8 and
531.0\,s. Extensive SAA is an exception to the certification
pattern for expected loss: it rises from 12/15 to 15/15 and 15/15,
while mean observed time falls from 578.0 to 401.8 and 249.5\,s.
These associations are descriptive, and times for unresolved runs
remain censored at the stopping rule.

These comparisons motivate Combinatorial B\&B as the preferred
expected-loss configuration and Path-exp as the preferred CVaR and
mean--CVaR configurations for a subsequent decision study. This
selection is based on certification and final gap, rather than an
unobserved out-of-sample outcome.

\FloatBarrier
\subsection{Out-of-sample stability of SAA solutions}\label{sec:oos-stability}

Does increasing $n$ improve both the average held-out objective and
its variation across the five sampled solutions? The appropriate
estimand evaluates one selected firebreak set $Y$ on the \emph{same}
1,000 reduced-graph scenarios for every $n$. For losses $q_s(Y)$ on
test IDs $s=9001,\ldots,10000$, define
\[
 \widehat\mu(Y)=\frac{1}{1000}\sum_s q_s(Y),\qquad
 \widehat C_{0.9}(Y)=\min_{z\in\mathbb R}
 \left\{z+\frac{1}{1000(1-0.9)}\sum_s(q_s(Y)-z)_+\right\}.
\]
The three test objectives are respectively $\widehat\mu$,
$\widehat C_{0.9}$, and $(1-0.5)\widehat\mu+0.5\widehat C_{0.9}$,
matching the optimization settings. Variation across the five
replications at each $(n,\alpha)$ should be reported, for example as
sample standard deviation, alongside the mean. Decision stability
requires the ten pairwise Jaccard similarities
$|Y_a\cap Y_b|/|Y_a\cup Y_b|$ within each five-replication group.
The training splits are independently generated at different sizes,
not nested: for example, \texttt{case00} at $\alpha=0.01$ shares only
four IDs between its 100- and 200-scenario training sets. Cross-size
\texttt{case\_id} therefore does not establish a paired design.

Combinatorial B\&B would be the preselected solution method for
expected loss, and Path-exp for CVaR and mean--CVaR, based solely on
Section~\ref{sec:results}. Using Path-exp across all objectives would
reduce expected-loss certification from 43/45 to 39/45, without
resolving its 24 uncertified CVaR or 15 uncertified mean--CVaR
problems. No placements are substituted across methods in this
decision analysis. Before attributing variation in held-out values
to sampling, unresolved sampled problems would need comparable
optimization accuracy. The task-level gaps and availability of
certified alternatives are recorded in
\texttt{tables/new20\_selection\_audit.csv}; this audit does not
measure held-out performance.

The target consolidated file displaces or truncates the placement
and evaluation fields, and its original worker solutions and source
scenario graphs are unavailable. Neither the selected original cell
sets nor their 1,000 scenario-level test losses can be independently
verified. An aggregate CVaR field cannot supply all three test
objectives. Therefore the mean and sample standard deviation of the
OOS objective for each of the 27 objective--budget--size combinations,
and the ten pairwise Jaccard similarities of each five-solution group,
cannot be estimated from these records. We do not report a partial
five-replication table selected by certification status or infer
stability at $n=400$.

\FloatBarrier
\subsection{Fidelity of the reduced propagation graph}\label{sec:graph-fidelity}

How much protection would the \emph{same} selected reduced-graph
placement retain on the full Reburn representation? For a fixed
placement $Y$ and the common test IDs, let $F_G(Y)$ be the empirical
risk functional above, evaluated using graph $G\in\{D,R\}$, where
$D$ is the reduced graph and $R$ the corresponding Reburn graph.
The paired raw transfer difference is $F_R(Y)-F_D(Y)$, and its relative
version is $[F_R(Y)-F_D(Y)]/F_D(Y)$ when $F_D(Y)>0$. When both
no-treatment baselines are available and nonzero, compare normalized
protection
\[
 P_G(Y)=\frac{F_G(\varnothing)-F_G(Y)}{F_G(\varnothing)},
 \qquad \Delta P(Y)=P_D(Y)-P_R(Y).
\]
These quantities compare fixed decisions, with pairs indexed by
$(\alpha,n,\text{replication},\text{objective},s,Y)$; they require no
Reburn optimization. The normalized difference accounts for distinct
no-treatment fire severity on the two graph representations.

The manifest verifies common test ID lists for reduced and Reburn
cases, but it cannot establish identical cell mappings, ignition
realizations, weights, or arc inclusion. The source scenario graphs
and the original solution files are absent. No scenario-level paired
loss or no-treatment baseline for these exact target pairs is
recoverable. Hence neither the transfer difference nor $\Delta P$
can be calculated. Proposition~\ref{prop:reduction} is conditional
on its graph-pair premises and is not used to assert empirical
monotonicity here. The independently optimized Reburn records do not
answer the fixed-placement question.
